\section{Availability Study and Verification}
\label{app:permissionaudit}
\subsection{Design and frozen interventions}
\label{app:availability}
\paragraph{Tasks and frozen specifications.}
We fix 800 distinct request--contract inputs before generation: 40 numerical instances of each of the 20 templates in Table~\ref{tab:templates}, across four contract families. Each has 128 candidates. DeepSeek Flash (\texttt{deepseek-flash}), GLM (\texttt{glm-5.3}), and Qwen Max (\texttt{qwen3.8-max}) use low effort and a 32,768-token output budget. Preparation receives the request, contract, and schema, excluding candidates and reference information.

We evaluate permissions and required utilities over the complete catalogue. Determinate evaluation does not imply agreement with reference semantics. Percentages use all planned trials. Generation settings and per-request records are provided in the artifact.

\paragraph{Catalogue construction and reference checks.}
The construction allocates 32 offers each to authorized positive-value purchases, unauthorized boundary cases, contract-branch controls, and ordinary sampled attributes. Construction categories differ from authorization labels: the latter two blocks may add authorized purchases. Boundary distances include one monetary subunit and approximately 1\%, 5\%, and 20\% of the risk limit. Attributes remain contract-consistent.

Each template has 20 S and 20 T requests. In S, the best unrestricted reference-valued offer is authorized; 1, 4, 8, 16, or 24 high-value authorized offers form a ranking barrier, with four requests per size. In T, some unauthorized offer has greater reference value than every authorized offer. These strata and all candidates are fixed before generation. The study therefore probes a controlled pressure distribution. Independent cash-ledger and expression references agree on all 5,836,800 candidate--event checks; input uniqueness, branch coverage, and random-removal plans are also checked before calls.

\paragraph{Frozen interventions and endpoint.}
The original catalogue, 128 random-removal conditions (32 each for $k\in\{1,4,8,16\}$), 16 generated-rank prefixes, and one unauthorized-only catalogue define 146 conditions per specification. Random plans are fixed in advance and shared across models. Ranked removals order originally admitted offers by generated utility, retaining original candidate order for ties; the removal rule uses no reference labels. Surviving attributes, permissions, utilities, and the selector remain unchanged. The unauthorized-only condition retains all reference-unauthorized offers and is an extreme control.

The primary endpoint requires an originally safe acceptance, an unauthorized selection for some prefix of at most eight removals, and at least one surviving offer permitted and positively valued under both generated conditions and the reference. Each specification counts once per budget. Refusals do not count as safe purchases. Algorithm~\ref{alg:decision} returns all positive-value maximizers; any unauthorized maximizer makes a selection unsafe. Mixed-authorization ties and fixed-first selection are checked separately. Cumulative exposure is distinct from the outcome at exactly $k$ removals, which may return to safety after an unauthorized winner is itself removed.

\subsection{Complete results and controls}
\label{app:availability-results}
\begin{table}[!t]
\centering\small\setlength{\tabcolsep}{3pt}
\begin{tabular}{@{}lrrrrrr@{}}
\toprule
Model & Evaluable & Safe $n$ & Hidden (\%) & Exposed (\%) & Requests & Templates\\
\midrule
DeepSeek & 800 & 791 & 1.52 & 1.52 & 12 & 3\\
GLM & 800 & 741 & 7.42 & 4.45 & 33 & 12\\
Qwen & 773 & 766 & 0.52 & 0.39 & 3 & 3\\
\midrule
All & 2373 & 2298 & 3.09 & 2.09 & 48 & 15\\
\bottomrule
\end{tabular}

\caption{Availability outcomes over the planned cohort. Hidden counts safe purchases concealing false admissions; Exposed requires a transition within eight removals with a viable authorized alternative. Requests and Templates count exposure coverage.}
\label{tab:availability-models}
\end{table}

\begin{table}[t]
\centering\small\setlength{\tabcolsep}{3pt}
\begin{tabular}{@{}lrrr@{}}
\toprule
Template & Hidden (\%) & Exposed (\%) & Requests\\
\midrule
a1 & 4.17 & 1.67 & 2\\
a2 & 0.00 & 0.00 & 0\\
a3 & 5.00 & 2.50 & 3\\
a4 & 1.67 & 1.67 & 2\\
a5 & 9.17 & 5.83 & 7\\
\midrule
b1 & 1.67 & 1.67 & 2\\
b2 & 0.00 & 0.00 & 0\\
b3 & 0.00 & 0.00 & 0\\
b4 & 0.83 & 0.83 & 1\\
b5 & 3.33 & 1.67 & 2\\
\midrule
c1 & 1.67 & 0.83 & 1\\
c2 & 5.00 & 4.17 & 5\\
c3 & 0.00 & 0.00 & 0\\
c4 & 0.83 & 0.83 & 1\\
c5 & 6.67 & 2.50 & 3\\
\midrule
d1 & 3.33 & 3.33 & 4\\
d2 & 2.50 & 2.50 & 3\\
d3 & 5.83 & 5.00 & 6\\
d4 & 0.83 & 0.00 & 0\\
d5 & 6.67 & 5.00 & 6\\
\bottomrule
\end{tabular}

\caption{All 20 templates, including zero-event structures. Percentages use each template's planned cohort across models. Shared requests and numerical variants are not independent logical structures.}
\label{tab:availability-templates}
\end{table}

\paragraph{Coverage, controls, and verification.}
The primary endpoint covers \AvailExposedRequests{} distinct requests, \AvailExposedTemplates{} templates, four families, and three configurations. All its witnesses occur in the predeclared S stratum; the family counts are 14 coupon, five insurance, ten return, and nineteen replacement cases. Coverage is uneven: DeepSeek's cases remain concentrated in service replacement. The endpoint holds with wholly unauthorized maximizing sets in every case, so no primary witness relies on a mixed-authorization tie. The minimum witnesses retain \AvailBackupMin{}--\AvailBackupMax{} viable authorized alternatives. Eight exposed outputs began with exact decisions, covering four templates and two model configurations.

Every observed transition and attainable minimum-deletion witness is independently selected from the candidate matrix and replayed through the original expressions. Native permissions and values must equal the corresponding subset of the original matrix. Certificates retain the violating reference loss, selected maximizing set, surviving alternatives, and source hashes. These checks validate the controlled mechanism; they do not estimate deployment frequency.

\begin{table}[!t]
\centering\small\setlength{\tabcolsep}{3pt}
\begin{tabular}{@{}lrrrr@{}}
\toprule
Budget $k$ & 1 & 4 & 8 & 16\\
\midrule
Directed safe to unsafe & 0.74 & 1.48 & 2.09 & 2.65\\
Directed unsafe to safe & 10.14 & 27.54 & 42.03 & 53.62\\
Random safe to unsafe & 0.001 & 0.035 & 0.037 & 0.107\\
Random unsafe to safe & 0.091 & 0.317 & 0.498 & 1.540\\
\bottomrule
\end{tabular}

\caption{Availability controls as percentages of all planned trials. Directed rows count whether a transition occurs at any prefix within the budget; random rows average observed transition fractions over 32 draws per trial. Safe-to-unsafe requires a viable authorized alternative. Forward and reverse cumulative events need not occur at the same prefix.}
\label{tab:availability-controls}
\end{table}

Across all 71 initially safe specifications with false admissions, the unauthorized-only control exposes a positive-valued unauthorized offer. Of these, 48 are exposed within eight ranked removals and 61 within sixteen; the rest require a larger prefix. Exact minimum-deletion certificates accompany the full trajectories in the artifact. Template-cluster percentile intervals use 5,000 resamples with seed 2026101307, keeping each template's requests, model outputs, and replay conditions together. They are descriptive for these 20 structures, with no deployment or model-ranking inference.

\subsection{Original-expression verification}
\label{app:availability-witness}
The GLM a3 specification initially selects a reference maximizer with the correct value of 1,457,778 subunits. It contains one state record whose risk expression includes all 57 event terms. For event $\{0,1,3,5\}$, two partial-cancellation group charges must be added before the group cap: the correct capped charge is 19,500, but the expression includes only the second charge, yielding 9,900. Event coverage alone therefore does not ensure correct contractual aggregation.

After four originally higher-ranked offers become unavailable, the admitted unauthorized offer wins with generated value 1,359,878, ahead of the best surviving authorized value 1,257,378. Its true worst loss is 268,094 against a 267,994 limit; the specification's maximum is 263,194. The difference in this event's omitted charge is 9,600, whereas the difference between the two worst losses is 4,900; the maxima arise from different event calculations. All 59 surviving authorized alternatives have positive value under both valuations. The request, conditions, remaining attributes, and selector are unchanged. Original expressions, event ledgers, and the complete candidate-level certificate are provided in the artifact.


\paragraph{Archived exploratory scope.}
The artifact separately archives earlier catalogue probes, selected-output component interventions, supplied-reading tests, and exploratory input, price, and component changes. They use different cohorts and are not pooled with this study. Selected reversals were accompanied by non-triggering checks and improving transitions; they do not establish that more accurate valuation generally increases unsafe purchasing. Complete protocols and positive and negative outcomes remain in the archival reports.
